% ===============================================
% MAIN FILE - ĐỀ CƯƠNG TOÁN 7
% ===============================================
\documentclass[14pt,a4paper]{extbook}
\usepackage[T5]{fontenc}
\usepackage[utf8]{inputenc}
\usepackage{lmodern} 

\makeatletter
\def\input@path{{../}{./}}
\makeatother

\usepackage{graphicx,multicol,tikz,amsmath,etoolbox,fontawesome5}
\usepackage{enumitem}
\usepackage{xcolor}
\usepackage{my-style}
\usepackage{setspace}

\onehalfspacing 
\let\Bbbk\relax
\let\openbox\relax

\begin{document}
\setmagiaotrinh{TOAN7-GHK1}
\pagestyle{sgkstyle}

\chude{ĐỀ CƯƠNG ÔN TẬP GHK1}

\begin{center}
    \textbf{\large NĂM HỌC: 2025-2026}
\end{center}

\phan{Lý thuyết}
\begin{dsSao}
    \item \textbf{Đại số:}
    \begin{dsTrangTri1}
        \item Tập hợp $\mathbb{Q}$ số hữu tỉ.
        \item Cộng, trừ, nhân, chia số hữu tỉ.
        \item Lũy thừa với số mũ tự nhiên của số hữu tỉ.
        \item Thứ tự thực hiện phép tính. Quy tắc chuyển vế.
        \item Biểu diễn thập phân của số hữu tỉ.
        \item Số vô tỉ, căn bậc hai số học.
    \end{dsTrangTri1}
    \item \textbf{Hình học:}
    \begin{dsTrangTri1}
        \item Nội dung chương III.
        \item Định lý tổng ba góc của một tam giác.
        \item Các trường hợp bằng nhau của tam giác.
    \end{dsTrangTri1}
\end{dsSao}

\phan{Bài tập}
\muc{BÀI TẬP TRẮC NGHIỆM}
\textit{Hãy khoanh tròn chữ cái in hoa đứng trước phương án đúng duy nhất}

\cau{Cách viết nào sau đây biểu diễn một số hữu tỉ?}
\begin{tracnghiem4}
    \item $\dfrac{4,5}{-5}$
    \item $\dfrac{3}{17}$
    \item $\dfrac{-100}{0}$
    \item $\dfrac{15}{0,4}$
\end{tracnghiem4}

\cau{Trong các khẳng định sau, khẳng định nào sai?}
\begin{tracnghiem4}
    \item $(-x)^{1}=-x$
    \item $x^{0}=1\ (x\ne0)$
    \item $x^{1}=1$
    \item $x^{a+1}=x^{a}\cdot x$
\end{tracnghiem4}

\cau{Số thập phân $0,36363636\ldots$ được viết dưới dạng thu gọn (có chu kì trong dấu ngoặc) là}
\begin{tracnghiem4}
    \item $0,(363)$
    \item $0,(36)$
    \item $0,(3)$
    \item $0,3(63)$
\end{tracnghiem4}

\muc{BÀI TẬP TỰ LUẬN}
\cau{Thực hiện các phép tính sau (một cách hợp lí nếu có thể):}
\begin{tuluan2}
    \item $9\cdot\left(-\dfrac{1}{3}\right)^{3}+\dfrac{1}{3}$
    \item $\dfrac{3}{7}\cdot 19\dfrac{1}{3}-\dfrac{3}{7}\cdot 33\dfrac{1}{3}$
    \item $\left[\dfrac{5}{8}+\left(-\dfrac{3}{4}\right)\right]+\dfrac{15}{6}$
    \item $15\dfrac{1}{4}:\left(-\dfrac{5}{7}\right)-25\dfrac{1}{4}:\left(-\dfrac{5}{7}\right)$
\end{tuluan2}

\end{document}